\section{Introduction}
\label{sec:introduction}

Wall-bounded turbulent flows occur in many engineering systems in which skin
friction, heat transfer, and mixing affect performance. Their behaviour is
governed by the interaction of solid boundaries, mean shear, and turbulent
motions over a broad range of scales. Pressure-driven plane channel flow is a
canonical configuration because its simple geometry isolates these mechanisms
and provides well-documented statistics for comparison
\cite{pope2000,kimMoinMoser1987}.

Direct numerical simulation (DNS) resolves the dynamically relevant turbulent
scales while avoiding a turbulence-closure model. It provides instantaneous
flow fields and converged statistics that support physical analysis, solver
assessment, and comparison with established benchmark data. Although DNS is
computationally demanding, plane channel flow remains a practical reference
case for evaluating mean and second-order turbulence statistics
\cite{kimMoinMoser1987}.

Temperature introduces an additional coupling when density variations produce
a buoyancy force in the momentum balance. The thermal field then becomes
active and can modify the velocity field. Under stable thermal stratification,
buoyancy resists wall-normal fluid displacement and can reduce turbulent
transport and mixing across the channel. Quantifying this interaction is
important for predicting momentum and heat transfer in stratified flows
\cite{armenioSarkar2002}.

The objective of the present project is to establish a reliable workflow for
studying these effects with the BRINE solver. The planned work first seeks to
establish or reproduce a non-buoyant turbulent channel-flow baseline and to
understand the solver setup, execution, and output-processing workflow. Active
thermal buoyancy will then be introduced to study how stable stratification
changes turbulence and wall-normal mixing. These are project objectives; their
outcomes will be assessed from the simulations and validation results reported
later.

The analysis will compare a non-buoyant case with stably stratified cases.
Where the confirmed solver configuration permits, the cases will use matched
operating conditions. The comparison will include the mean streamwise velocity,
RMS velocity fluctuations, Reynolds shear stress, and turbulent kinetic energy,
together with thermal quantities defined in later work. Agreement of the
non-buoyant baseline with trusted DNS data is required before differences in
the buoyant cases are interpreted.

The remainder of the thesis is organised around this progression. The
theoretical background first defines the channel configuration, governing
equations, and baseline turbulence statistics. Subsequent chapters will
document the BRINE workflow and numerical configuration, assess the
non-buoyant baseline, introduce active thermal buoyancy, and compare the
stratified and non-stratified results. The final chapter will summarise the
main findings and identify remaining work.
